\documentclass[10pt,a4paper]{amsart}
\usepackage[margin=19mm]{geometry}
\usepackage{amsmath,amssymb,mathtools}
\usepackage[scr=rsfs,cal=euler]{mathalfa}
\usepackage{xfrac}
\usepackage[unicode,hidelinks,bookmarks=false]{hyperref}
\newcommand{\ie}{i.e.\ }
\newcommand{\cf}{cf.\ }
\newcommand{\resp}{respectively\ }
\newcommand{\Subsection}{\subsection}
\providecommand{\nobreakdash}{\nobreak\hbox{-}}
\setlength{\emergencystretch}{2em}
\multlinegap0pt
\usepackage{bm}
\usepackage{bbm}
\usepackage{dsfont}
\usepackage{xspace}
\usepackage{cancel}
\usepackage{mathtools}
\usepackage{amsthm}
\makeatletter
\@ifundefined{theorem}{\newtheorem{theorem}{Theorem}}{}
\makeatother
\title[Existence Result for Singular Second Order Dynamic Equations]{Existence Result for Singular Second Order Dynamic Equations with Mixed Boundary Conditions}
\author{Shalmali Bandyopadhyay, Curtis J Kunkel}
\date{Source: arXiv:2506.16505v1 (CC BY 4.0)}
\begin{document}
\maketitle

\noindent\emph{Source and licence.} Existence Result for Singular Second Order Dynamic Equations with Mixed Boundary Conditions. Version arXiv:2506.16505v1 (CC BY 4.0), distributed under the Creative Commons Attribution 4.0 International licence, as recorded in the arXiv licence metadata for this exact version and shown on the abstract page https://arxiv.org/abs/2506.16505v1. Full rights notice: licenses/arxiv-2506.16505-CC-BY-4.0.txt.\par\smallskip

\noindent\textit{Editorial scope.} Counted unit: the Theorem whose source label is \texttt{main} in the article (source0.tex lines 337--346, source label \texttt{main}), delivered together with the article's own complete original proof of it (source0.tex lines 347--405, one \texttt{proof} environment of 59 lines). No mathematics is written, repaired or re-derived: statement and proof are the article's own text line for line. Required context retained uncounted: eq (lines 108--114); eq\_reg (lines 191--197); thm\_reg (lines 228--230). Every definition, display and label of those lines is retained unchanged. Omitted from this package and not redistributed: 1--107, 115--190, 198--227, 231--336, 406--512. Nothing inside a retained excerpt is omitted or altered: every retained body is delivered exactly as the article's lines are written, so no omission receipt of any kind accompanies this package. Citations: the retained excerpts carry no citation command at all, so no bibliography entry of the article is transcribed and no reference is redistributed. Reference adaptation: none was needed and none was made; every reference inside the delivered unit resolves inside it, so no unresolved reference or citation is produced. Printed slips kept literal: no printed word, symbol, hypothesis or number of the counted statement or of its proof was altered. Layout and class: the article's own class is \texttt{amsart} with packages including amssymb, amsthm, latexsym, amsmath, url, color, fullpage, hyperref; the wrapper is the lane's pinned \texttt{amsart} editorial wrapper at 10pt on A4, which restates the theorem-like environments the retained text needs and adds the article's own macro definitions that the retained text needs (none), with extra wrapper packages (bm, bbm, dsfont, xspace, cancel, mathtools). The delivered numbering is the unit's own. Assets: no retained span contains a figure environment, a graphics-inclusion command, a PDF-page inclusion, a table or a TikZ picture; no image file is copied into this package, the package declares no asset dependency and no image file is redistributed with it. Licence: version arXiv:2506.16505v1 of this article is distributed under the Creative Commons Attribution 4.0 International licence, as recorded in the arXiv licence metadata for this exact version and shown on the abstract page https://arxiv.org/abs/2506.16505v1. A full rights notice is delivered at licenses/arxiv-2506.16505-CC-BY-4.0.txt. Characters: every retained excerpt is pure ASCII, so the delivered engine, XeLaTeX with its default Latin Modern text font, reports no missing character.
\begin{equation}\label{eq}
\begin{cases}
u^{\Delta\Delta}(\rho(t)) + f(t, u(t)) = 0, \quad t \in \mathbb{T}^0,\\
u^{\Delta}(0) = 0, \\
u(T) = g(T).
\end{cases}
\end{equation}

\begin{equation}\label{eq_reg}
\begin{cases}
u^{\Delta\Delta}(u(t)) + h(t, u(t)) = 0, \quad t \in \mathbb{T}^0\\
u^{\Delta}(0)=0\\
u(T)=g(T)
\end{cases}
\end{equation}

\begin{theorem}\label{thm_reg}
	Let $\alpha$ and $\beta$ be lower and upper solutions of the regular problem \eqref{eq_reg} respectively, such that, $\alpha \leq \beta$ on $\mathbb{T}.$  Let $h(t, x, y)$ be continuous on $\mathbb{T} \times \mathbb{R}^2$ and non-increasing in its $y$-variable.  Then \eqref{eq_reg} has a solution $u$ satisfying $$\alpha(t) \leq u(t) \leq \beta(t), \quad t \in \mathbb{T}.$$
\end{theorem}

\begin{theorem}\label{main}
	Assume conditions \emph{(\textbf{A})}, \emph{(\textbf{B})}, and \emph{(\textbf{C})} hold, along with the following:
	\begin{description}
		\item[D] There exists $c \in (0,\infty)$ so that $f(t,c) \leq 0,$ for all $t \in \mathbb{T}^0.$
		\item[E] There exists $\delta > 0$ so that $f(t,x) > 0$ for all $t \in (T-\delta,T) \cap \mathbb{T}$ and $x \in \left(0, \frac{c}{2}\right).$
		\item[F] $\lim\limits_{x \rightarrow 0^{+}} f(t,x) = \infty$ for $t \in \mathbb{T}.$ 
        \item[G] $0< g(T) <u(0)$ 
	\end{description}
    Then, \eqref{eq} has a solution $u$ satisfying $$0 < u(t) \leq c, \quad t \in \mathbb{T}^k.$$
\end{theorem}

\begin{proof}
	We begin by modifying our given function $f$ as follows.
	
	For $k > 0, t \in \mathbb{T},$ define
	\begin{equation*}
	f_{k}(t,x) = \left\{\begin{array}{ll}
	f\left(t,|x|\right), & \text{if } |x| \geq \frac{1}{k} \\
	f\left(t,\frac{1}{k}\right), & \text{if } |x| < \frac{1}{k}.
	\end{array}\right.
	\end{equation*}
	Then, $f_{k}$ is continuous on $\mathbb{T} \times \mathbb{R}.$
	
	Assumption (\textbf{F}) implies that there exists a $k_{0}$ such that, for all $k \geq k_{0},$ $$f_{k}(t,0) = f\left(t,\frac{1}{k}\right) > 0, \quad \text{for all } t \in \mathbb{T}.$$
	
	We now consider the boundary value problem
	\begin{equation}\label{eq_k}
	u^{\Delta\Delta}(\rho(t)) + f_{k}(t, u(t)) = 0, \quad t \in \mathbb{T}^0,
	\end{equation}
	satisfying boundary conditions \eqref{eq}.
	
	Now, let $\alpha(t) = 0$ and $\beta(t) = c.$  Then, for each $k \geq k_{0},$ $\alpha$ and $\beta$ are lower and upper solutions of \eqref{eq_k} respectively.  Also, $\alpha(t) \leq \beta(t)$ for $t \in \mathbb{T}.$  Thus, by Theorem \ref{thm_reg}, for each $k \geq k_{0},$ there exists a solution $u_k$ to each problem \eqref{eq_k} that satisfies $0 \leq u_{k}(t) \leq c,$ for $t \in \mathbb{T}.$
	
	Let $\mathbb{T}_1 = \mathbb{T} \cap (0, T-\delta)$ and $\mathbb{T}_2 = \mathbb{T} \cap (T-\delta, T).$
	
	Now, for each $k \geq k_0,$ there exists $\varepsilon \in \left(0, \frac{c}{2}\right)$ such that if $k_{\varepsilon} \geq k_{0},$ then 
	\begin{equation}\label{fk}
		f_{k}(t, x) > c, \quad t \in \mathbb{T}, x \in (0, \varepsilon].
	\end{equation}
	
	For the sake of establishing a contradiction, assume that for $k \geq k_{\varepsilon} \geq k_0,$ we have that $u_{k_{\varepsilon}}(t) < \varepsilon_1,$ where 
	\begin{equation*}
		\varepsilon_1 = \left\{
		\begin{array}{ll}
			\varepsilon, & t \in \mathbb{T}_1, \\
			\frac{\varepsilon}{\delta}(T-t), & t \in \mathbb{T}_2.
		\end{array}
		\right.
	\end{equation*}
	
	Now,
	\begin{eqnarray*}
		-u_{k_{\varepsilon}}(t) & = &-g(T)- \int_{t}^{T} u_{k_{\varepsilon}}^{\Delta}(r) \Delta r \\
		& = & -g(T)- \int_{t}^{T} \int_{0}^{r} f_{k_{\varepsilon}}(s, u_{k_{\varepsilon}}(s)) \Delta s \Delta r \\
		& = & -g(T)- \int_{t}^{T-\delta} \int_{0}^{r} f_{k_{\varepsilon}}(s, u_{k_{\varepsilon}}(s)) \Delta s \Delta r \\
		&   & - \int_{T-\delta}^{T} \int_{0}^{r} f_{k_{\varepsilon}}(s, u_{k_{\varepsilon}}(s)) \Delta s \Delta r \\
		& < & - \int_{T-\delta}^{T} \int_{0}^{r} c \Delta s \Delta r \\
		& = & - \frac{c}{2} \left((T-\delta)^2 - t^2\right).
	\end{eqnarray*}
	
	First, consider $t \in \mathbb{T}_1.$  We have that $$u_{k_{\varepsilon}}(t) > \frac{c}{2} \left((T-\delta)^2 - t^2\right).$$  However, this implies that $u_{k_{\varepsilon}}^{\Delta\Delta}(\rho(t)) > -c,$ which leads to $f_{k_{\varepsilon}}(t, u_{k_{\varepsilon}}(t)) < c,$ a contradiction to \eqref{fk}.  Hence, for $t \in \mathbb{T}_1,$ we have $$u_{k_{\varepsilon}}(t) \geq \varepsilon_1 = \varepsilon.$$
	
	Now, consider $t \in \mathbb{T}_2.$  We have that $u_{k_{\varepsilon}} > \frac{c}{2} \left((T-\delta)^2 - t^2\right).$  This again implies that $u_{k_{\varepsilon}}^{\Delta\Delta}(\rho(t)) > -c.$  We also have, from assumption (\textbf{E}), that $$u_{k_{\varepsilon}}^{\Delta\Delta}(\rho(t)) = - f_{k_{\varepsilon}}(t, u_{k_{\varepsilon}}(t)) < 0.$$  Hence, $$-c < u_{k_{\varepsilon}}^{\Delta\Delta}(\rho(t)) < 0,$$ making $u_{k_{\varepsilon}}$ concave in this interval.  We also know, by continuity, that $u_{k_{\varepsilon}}(T-\delta) \geq \varepsilon$ and $u_{k_{\varepsilon}}(T) = 0.$  Therefore, by continuity, for $t \in \mathbb{T}_2,$ $$u_{k_{\varepsilon}}(t) > \varepsilon_1 = \frac{\varepsilon}{\delta} (T-t).$$
	
	Hence, $0 < \varepsilon_1 < u_{k_{\varepsilon}}(t)$ for $t \in \mathbb{T}^0.$
	Thus, for $k \geq k_{\varepsilon},$ we can choose a subsequence $\left\{u_{k_{n}}(t)\right\} \subseteq \left\{u_{k}(t)\right\}$ so that $$\lim\limits_{n \rightarrow \infty} u_{k_{n}}(t) = u(t), \quad t \in \mathbb{T},$$ and note that $u(t) \in E,$ where $E$ is defined as in the proof of Theorem \ref{thm_reg}.
     
    Moreover, for sufficiently large $n,$ $$u_{k_{n}} = g(T)+\int_{t}^{T} \int_{0}^{r} f(s, u_{k_{n}}(s)) \Delta s \Delta r.$$  And from the continuity of $f,$ as we let $n \rightarrow \infty,$ we get $$u(t) =g(T)\int_{t}^{T} \int_{0}^{r} f(s, u(s)) \Delta s \Delta r.$$  Hence, $$u^{\Delta\Delta}(\rho(t)) = -f(t, u(t)),$$ with the desired inequality satisfied, specifically, $0 < u(t) < c$ on $\mathbb{T}^k.$

\end{proof}

\end{document}
